\section{Update protocol and state schema}
\label{sec:protocol}

The hybrid comparison crosses the update interface with the semantic
schema (Table~\ref{tab:protocol}). The tool-based reference with the existing
schema reaches \HybridRate\%. Replacing that schema with the task-adapted
version raises the observed solved rate to \TaskAdaptedRate\%, a
\TaskAdaptedVsHybridDelta{} percentage-point change. Inline updates with
the existing schema instead reach \InlineRate\%, a
$\InlineVsHybridDelta$ point change from the reference.

\begin{table}[htbp]
\centering\small
\caption{Update protocol and schema within hybrid state. Output tokens
and cost are per-execution medians. Exact McNemar tests compare each
configuration with Tool + Existing schema.}
\label{tab:protocol}
\input{tables/new-protocol}
\end{table}

Combining inline updates with the adapted schema yields
\InlineAdaptedRate\%, below either change in isolation. The descriptive
difference-in-differences for protocol and schema is
$\ProtocolSchemaInteraction$ percentage points. This is the difference
between the inline effect under the adapted schema and its effect under
the existing schema, not the joint configuration's difference from
the reference. The available aggregates do not support a task-level
interaction test.

Neither intervention recovers the within-comparison full-history rate of
\HybridControlRate\%. The result does not support the hypothesis that
these particular protocol or schema changes are sufficient to repair the
observed deficit. It also does not rule out better schemas or more faithful
implementations of the published runtime.

The inline representation changes both how actions are expressed and when
state maintenance occurs. Its larger reported output median may include
additional patch content or changes elsewhere in the trajectory.
Attributing the success difference to formatting burden, patch validity,
or a specific omitted field would require a controlled trace-level
analysis. Schema adaptation is likewise an implementation choice here,
not an independently validated memory abstraction.
